\subsection{Generalized Eigenspaces}

\begin{definition} \label{def:generalized-eigenvector-and-eigenspace} \hlt{Generalized Eigenvectors and Generalized Eigenspaces}\\
Let $T\in\LL(\mathcal{V})$, $\lambda\in\mathbb{F}$. Nonzero vector $\mathbf{v}\in\mathcal{V}$ is a \hlt{generalized eigenvector} of $T$ to $\lambda$ if $(T-\lambda I_{\mathcal{V}})^j\mathbf{v}=\mathbf{0}$ for some positive integer $j$. The \hlt{generalized eigenspace} corresponding to $\lambda$ is
\begin{align}
\mathcal{G}_{\lambda}(T)\defeq \{\mathbf{0}\}\cup\{\mathbf{v}\neq\mathbf{0}:\mathbf{v}\text{ is a generalized eigenvector of }T\text{ corresponding to }\lambda\}.
\nonumber
\end{align}
\end{definition}

Every eigenvector is a generalized eigenvector by taking $j=1$. The point is that higher powers of $T-\lambda I_{\mathcal{V}}$ may reveal directions that are not eigenvectors but are eventually pushed into the eigenspace.

\begin{proposition} \label{prop:kernels-of-powers-increase} \hlt{Kernels of Powers Increase}\\
Let $T\in\LL(\mathcal{V})$. Then $\{\mathbf{0}\}=\ker(T^0)\subseteq\ker(T)\subseteq\ker(T^2)\subseteq\cdots\subseteq\ker(T^k)\subseteq\ker(T^{k+1})\subseteq\cdots$.
\begin{proof}
If $\mathbf{v}\in\ker(T^k)$, then $T^k\mathbf{v}=\mathbf{0}$, hence $T^{k+1}\mathbf{v}=T(T^k\mathbf{v})=\mathbf{0}$. Thus $\ker(T^k)\subseteq\ker(T^{k+1})$ for every $k\geq0$.\\
\hlt{Structural Commentary}: Powers of $T$ form an increasing chain of kernels.\\
\hlt{Key Insight}: Once a vector is killed by $T^k$, applying one more $T$ cannot revive it.
\end{proof}
\end{proposition}

\begin{proposition} \label{prop:kernel-power-stabilization-after-equality} \hlt{Equality Forces Stabilization}\\
Let $T\in\LL(\mathcal{V})$ and $m\geq0$. If $\ker(T^m)=\ker(T^{m+1})$, then $\ker(T^m)=\ker(T^{m+1})=\ker(T^{m+2})=\cdots$.
\begin{proof}
By Proposition \ref{prop:kernels-of-powers-increase}, it suffices to prove $\ker(T^{m+k+1})\subseteq\ker(T^{m+k})$ for every $k\geq0$.\\
Take $\mathbf{v}\in\ker(T^{m+k+1})$. Then $T^{m+1}(T^k\mathbf{v})=T^{m+k+1}\mathbf{v}=\mathbf{0}$, so $T^k\mathbf{v}\in\ker(T^{m+1})=\ker(T^m)$.\\
Hence $T^{m+k}\mathbf{v}=T^m(T^k\mathbf{v})=\mathbf{0}$, proving the reverse inclusion.\\
\hlt{Structural Commentary}: Stabilization at one step propagates through the whole tail.\\
\hlt{Key Insight}: A later kernel vector can be shifted down into the already-stable kernel.
\end{proof}
\end{proposition}

\begin{theorem} \label{thm:kernel-powers-stabilize-by-dimension} \hlt{Kernels Stop Growing by $\dim(\mathcal{V})$}\\
If $\mathcal{V}$ is finite-dimensional, $T\in\LL(\mathcal{V})$, and $n=\dim(\mathcal{V})$, then $\ker(T^n)=\ker(T^{n+1})=\ker(T^{n+2})=\cdots$.
\begin{proof}
If $\ker(T^n)\neq\ker(T^{n+1})$, Prop \ref{prop:kernel-power-stabilization-after-equality} implies all inclusions $\{\mathbf{0}\}=\ker(T^0)\subsetneq\ker(T)\subsetneq\cdots\subsetneq\ker(T^{n+1})$
are strict. Each strict inclusion raises dimension by at least $1$, so $\dim\ker(T^{n+1})\geq n+1$, impossible because $\ker(T^{n+1})\subseteq\mathcal{V}$. Thus $\ker(T^n)=\ker(T^{n+1})$, and Proposition \ref{prop:kernel-power-stabilization-after-equality} gives the tail equality.\\
\hlt{Structural Commentary}: Strict kernel growth is dimension-limited.\\
\hlt{Key Insight}: Dimension bounds force the kernel chain to stabilize no later than step $\dim(\mathcal{V})$.
\end{proof}
\end{theorem}

\begin{figure}[H]
  \centering
  \begin{tikzpicture}[
  font=\small,
  subspace/.style={draw=blue!60, fill=blue!6, rounded corners=2pt, minimum height=0.65cm, align=center},
  arrow/.style={-{Stealth[length=2mm]}, thick, blue!65}
]
  \node[subspace, minimum width=1.2cm] (k0) at (0,0) {$\ker(T^0)$};
  \node[subspace, minimum width=1.5cm] (k1) at (2.1,0) {$\ker(T)$};
  \node[subspace, minimum width=1.7cm] (k2) at (4.4,0) {$\ker(T^2)$};
  \node[subspace, minimum width=1.2cm] (dots) at (6.5,0) {$\cdots$};
  \node[subspace, minimum width=1.7cm] (kn) at (8.4,0) {$\ker(T^n)$};
  \node[subspace, minimum width=1.8cm] (kn1) at (11.5,0) {$\ker(T^{n+1})$};
  \node[subspace, minimum width=1.2cm] (tail) at (13.5,0) {$\cdots$};

  \draw[arrow] (k0) -- (k1);
  \draw[arrow] (k1) -- (k2);
  \draw[arrow] (k2) -- (dots);
  \draw[arrow] (dots) -- (kn);
  \draw[arrow] (kn) -- node[above, black] {\scriptsize stabilizes} (kn1);
  \draw[arrow] (kn1) -- (tail);

  \draw[gray!70, thick] (-0.7,-0.75) -- (8.75,-0.75);
  \draw[gray!70, thick] (-0.7,-0.65) -- (-0.7,-0.85);
  \draw[gray!70, thick] (8.75,-0.65) -- (8.75,-0.85);
  \node[below, black] at (4.1,-0.95) {at most $\dim(\mathcal{V})$ strict growth steps};
\end{tikzpicture}

  \caption{The kernel filtration grows until it stabilizes; after equality appears once, the tail is constant.}
  \label{fig:kernel-power-filtration}
\end{figure}

\begin{theorem} \label{thm:ranks-of-matrix-powers-stabilize} \hlt{Ranks of Matrix Powers Stabilize}\\
Let $A\in\operatorname{Mat}(n,n,\mathbb{F})$. Then:
\begin{enumerate}[label=(\roman*)]
\setlength{\itemsep}{0pt}
\item $\operatorname{rank}(A^k)\geq\operatorname{rank}(A^{k+1})$ for $k=1,2,\ldots$;
\item if $\operatorname{rank}(A^k)=\operatorname{rank}(A^{k+1})$, then $\operatorname{col}(A^k)=\operatorname{col}(A^{k+p})$ and $\operatorname{rank}(A^k)=\operatorname{rank}(A^{k+p})$ for every $p=1,2,\ldots$;
\item if $A$ is not invertible, then there is a least positive integer $q\leq n$ such that $\operatorname{rank}(A^q)=\operatorname{rank}(A^{q+1})$.
\end{enumerate}
\begin{proof}
\begin{enumerate}[label=(\roman*)]
\setlength{\itemsep}{0pt}
\item Since $A^{k+1}=A^kA$, each column of $A^{k+1}$ is a linear combination of the columns of $A^k$. Hence $\operatorname{col}(A^{k+1})\subseteq\operatorname{col}(A^k)$, so $\operatorname{rank}(A^{k+1})\leq\operatorname{rank}(A^k)$.
\item The inclusion $\operatorname{col}(A^{k+1})\subseteq\operatorname{col}(A^k)$ and equality of dimensions give $\operatorname{col}(A^{k+1})=\operatorname{col}(A^k)$. Let $U=\operatorname{col}(A^k)$. Then multiplication by $A$ maps $U$ onto $U$, so induction gives $\operatorname{col}(A^{k+p})=U$ for every $p\geq1$.
\item If $A$ is not invertible, then Corollary \ref{cor:square-matrix-invertible-iff-full-rank} gives $\operatorname{rank}(A)\leq n-1$. If all inequalities $\operatorname{rank}(A^k)>\operatorname{rank}(A^{k+1})$ held for $k=1,\ldots,n$, then $\operatorname{rank}(A^{n+1})\leq\operatorname{rank}(A)-n<0$, impossible. Hence equality occurs for some $q\leq n$; choose the least such $q$.
\end{enumerate}
\hlt{Structural Commentary}: The column spaces of powers form a decreasing chain. In finite dimensions, a decreasing chain cannot keep losing dimension forever.\\
\hlt{Key Insight}: Once two consecutive ranks agree, the image has stabilized under multiplication by $A$.
\end{proof}
\end{theorem}

\begin{definition} \label{def:matrix-index} \hlt{Index of a Matrix}\\
If $A\in\operatorname{Mat}(n,n,\mathbb{F})$ is invertible, define its \hlt{index} to be $0$. If $A$ is not invertible, its index is the least positive integer $q$ such that $\operatorname{rank}(A^q)=\operatorname{rank}(A^{q+1})$.
\end{definition}

\begin{theorem} \label{thm:kernel-range-stable-direct-sum} \hlt{Stable Kernel-Range Decomposition}\\
Let $\mathcal{V}$ be finite-dimensional, let $T\in\LL(\mathcal{V})$, and set $n=\dim(\mathcal{V})$. Then $\mathcal{V}=\ker(T^n)\oplus\operatorname{range}(T^n)$.
\begin{proof}
First show the intersection is trivial. If $\mathbf{v}\in\ker(T^n)\cap\operatorname{range}(T^n)$, then $\mathbf{v}=T^n\mathbf{u}$ and $T^n\mathbf{v}=\mathbf{0}$. Hence $T^{2n}\mathbf{u}=\mathbf{0}$, so $\mathbf{u}\in\ker(T^{2n})=\ker(T^n)$ by Theorem \ref{thm:kernel-powers-stabilize-by-dimension}. Thus $\mathbf{v}=T^n\mathbf{u}=\mathbf{0}$.\\
By Rank-Nullity Theorem \ref{thm:rank-nullity}, $\dim\ker(T^n)+\dim\operatorname{range}(T^n)=\dim(\mathcal{V})$.\\
The zero intersection and the dimension count imply the sum is direct and equals $\mathcal{V}$.\\
\hlt{Structural Commentary}: At the stable power, killed directions and surviving image directions separate cleanly.\\
\hlt{Key Insight}: Stabilization turns the usual rank-nullity count into an actual direct-sum decomposition.
\end{proof}
\end{theorem}

\begin{theorem} \label{thm:generalized-eigenspace-stabilized-kernel} \hlt{Description of Generalized Eigenspaces}\\
Let $\mathcal{V}$ be finite-dimensional, $T\in\LL(\mathcal{V})$, $\lambda\in\mathbb{F}$, and $n=\dim(\mathcal{V})$. Then $\mathcal{G}_{\lambda}(T)=\ker\bigl((T-\lambda I_{\mathcal{V}})^n\bigr)$.\\
In particular, $\mathcal{G}_{\lambda}(T)$ is a subspace of $\mathcal{V}$.
\begin{proof}
Set $S=T-\lambda I_{\mathcal{V}}$.
\begin{enumerate}[label=\roman*.]
\setlength{\itemsep}{0pt}
\item[$(\Rightarrow)$] If $\mathbf{v}\in \mathcal{G}_{\lambda}(T)$, then $S^j\mathbf{v}=\mathbf{0}$ for some $j\geq1$. If $j\leq n$, Proposition \ref{prop:kernels-of-powers-increase} gives $\mathbf{v}\in\ker(S^n)$; if $j>n$, Theorem \ref{thm:kernel-powers-stabilize-by-dimension} gives $\ker(S^j)=\ker(S^n)$.
\item[$(\Leftarrow)$] If $\mathbf{v}\in\ker(S^n)$, then either $\mathbf{v}=\mathbf{0}$ or $\mathbf{v}$ is a generalized eigenvector of $T$ corresponding to $\lambda$.
\end{enumerate}
Thus $\mathcal{G}_{\lambda}(T)=\ker(S^n)$, which is a subspace.\\
\hlt{Structural Commentary}: Generalized eigenspaces are stable shifted kernels.\\
\hlt{Key Insight}: Finite dimension makes exponent $\dim(\mathcal{V})$ enough.
\end{proof}
\end{theorem}

\begin{proposition} \label{prop:generalized-eigenspaces-invariant} \hlt{Generalized Eigenspaces Are Invariant}\\
Let $\mathcal{V}$ be finite-dimensional and $T\in\LL(\mathcal{V})$. Then $\mathcal{G}_{\lambda}(T)$ is invariant under $T$ for every $\lambda\in\mathbb{F}$.
\begin{proof}
By Theorem \ref{thm:generalized-eigenspace-stabilized-kernel}, $\mathcal{G}_{\lambda}(T)=\ker((T-\lambda I_{\mathcal{V}})^n)$, where $n=\dim(\mathcal{V})$.\\
This is the kernel of a polynomial in $T$, so it is invariant under $T$ by Proposition \ref{prop:kernel-of-polynomial-in-operator-is-invariant}.\\
\hlt{Structural Commentary}: A generalized eigenspace is cut out by a polynomial equation in $T$.\\
\hlt{Key Insight}: Polynomial kernels commute with the operator that defines them.
\end{proof}
\end{proposition}

\begin{definition} \label{def:nilpotent-operator} \hlt{Nilpotent Operator}\\
An operator $N\in\LL(\mathcal{V})$ is \hlt{nilpotent} if $N^m=0$ for some positive integer $m$.
\end{definition}

\begin{corollary} \label{cor:nilpotent-power-dimension-zero} \hlt{Nilpotent Power Bound}\\
If $\mathcal{V}$ is finite-dimensional, $N\in\LL(\mathcal{V})$ is nilpotent, and $n=\dim(\mathcal{V})$, then $N^n=0$.
\begin{proof}
Since $N$ is nilpotent, $\ker(N^m)=\mathcal{V}$ for some $m$.\\
The kernel chain has stabilized by the $n$th step by Theorem \ref{thm:kernel-powers-stabilize-by-dimension}, so $\ker(N^n)=\mathcal{V}$. Hence $N^n=0$.\\
\hlt{Structural Commentary}: Nilpotence says some kernel in the chain is already the whole space.\\
\hlt{Key Insight}: The stabilization bound pulls that full-kernel conclusion back to exponent $\dim(\mathcal{V})$.
\end{proof}
\end{corollary}

\begin{lemma} \label{lem:scalar-plus-nilpotent-invertible} \hlt{Nonzero Scalar Plus Nilpotent Is Invertible}\\
Let $N\in\LL(\mathcal{V})$ be nilpotent and let $a\in\mathbb{F}$ with $a\neq0$. Then $aI_{\mathcal{V}}+N$ is invertible.
\begin{proof}
Choose $m$ such that $N^m=0$. Then $(aI_{\mathcal{V}}+N)\left(a^{-1}\sum\limits_{r=0}^{m-1}(-a^{-1}N)^r\right) =I_{\mathcal{V}}-(-a^{-1}N)^m=I_{\mathcal{V}}$,
and the same computation in the reverse order gives a two-sided inverse.\\
\hlt{Structural Commentary}: A nilpotent perturbation has a finite geometric-series inverse.\\
\hlt{Key Insight}: Nilpotence truncates the Neumann series exactly.
\end{proof}
\end{lemma}

\begin{theorem} \label{thm:distinct-generalized-eigenvectors-independent} \hlt{Generalized Eigenvectors for Distinct Eigenvalues Are Independent}\\
Let $\mathcal{V}$ be finite-dimensional and $T\in\LL(\mathcal{V})$. Suppose $\lambda_1,\ldots,\lambda_m$ are distinct eigenvalues of $T$, and let $\mathbf{v}_1,\ldots,\mathbf{v}_m$ be corresponding generalized eigenvectors. Then $(\mathbf{v}_1,\ldots,\mathbf{v}_m)$ is linearly independent.
\begin{proof}
Let $n=\dim(\mathcal{V})$ and suppose $a_1\mathbf{v}_1+\cdots+a_m\mathbf{v}_m=\mathbf{0}$.\\
Fix $r\in\{1,\ldots,m\}$ and apply $Q_r(T)\defeq\prod\limits_{i\neq r}(T-\lambda_iI_{\mathcal{V}})^n$ to the relation.\\
For $i\neq r$, Theorem \ref{thm:generalized-eigenspace-stabilized-kernel} gives $Q_r(T)\mathbf{v}_i=\mathbf{0}$, so $a_rQ_r(T)\mathbf{v}_r=\mathbf{0}$.\\
On $\mathcal{G}_{\lambda_r}(T)$, the operator $T-\lambda_iI_{\mathcal{V}}$ equals $(T-\lambda_rI_{\mathcal{V}})+(\lambda_r-\lambda_i)I_{\mathcal{V}}$. The first term is nilpotent on $\mathcal{G}_{\lambda_r}(T)$ and $\lambda_r-\lambda_i\neq0$, so Lemma \ref{lem:scalar-plus-nilpotent-invertible} makes each factor invertible on $\mathcal{G}_{\lambda_r}(T)$. Thus $Q_r(T)\mathbf{v}_r\neq\mathbf{0}$, forcing $a_r=0$. Since $r$ was arbitrary, all coefficients vanish.\\
\hlt{Structural Commentary}: Polynomial filters isolate one generalized eigenspace at a time.\\
\hlt{Key Insight}: Distinct eigenvalues let polynomial filters kill the other components and preserve the target one.
\end{proof}
\end{theorem}

\begin{theorem} \label{thm:nilpotent-strictly-upper-triangular-matrix} \hlt{Nilpotent Operators and Strictly Upper-Triangular Matrices}\\
Let $\mathcal{V}$ be finite-dimensional and $N\in\LL(\mathcal{V})$. Then $N$ is nilpotent if and only if $N$ has a strictly upper-triangular matrix in some basis.
\begin{proof}
Let $n=\dim(\mathcal{V})$. By Corollary \ref{cor:nilpotent-power-dimension-zero}, $\{\mathbf{0}\}\subseteq\ker(N)\subseteq\ker(N^2)\subseteq\cdots\subseteq\ker(N^n)=\mathcal{V}$.
\begin{enumerate}[label=\roman*.]
\setlength{\itemsep}{0pt}
\item[$(\Rightarrow)$] Choose a basis of $\ker(N)$, extend it to a basis of $\ker(N^2)$, and continue until obtaining a basis of $\mathcal{V}$.\\
If $\mathbf{v}$ first appears in $\ker(N^k)$, then $N\mathbf{v}\in\ker(N^{k-1})$, which is spanned by earlier basis vectors. Hence the matrix of $N$ is strictly upper triangular.
\item[$(\Leftarrow)$] If the matrix of $N$ is $n\times n$ strictly upper triangular, then $N$ sends each basis vector into the span of earlier basis vectors. After $n$ applications every basis vector is sent to $\mathbf{0}$, so $N^n=0$.
\end{enumerate}
\hlt{Structural Commentary}: Kernel filtration gives adapted basis; triangular form records motion to earlier layers.\\
\hlt{Key Insight}: Nilpotence is strict upper-triangularity in a suitable basis.
\end{proof}
\end{theorem}

\begin{figure}[H]
  \centering
  \begin{tikzpicture}[
  font=\small,
  layer/.style={draw=teal!65, fill=teal!7, rounded corners=2pt, minimum width=1.95cm, minimum height=0.55cm, align=center},
  arrow/.style={-{Stealth[length=2mm]}, thick, teal!70}
]
  \node[layer] (n1) at (0,0) {$\ker(N)$};
  \node[layer] (n2) at (2.46,0) {$\ker(N^2)$};
  \node[layer] (n3) at (4.93,0) {$\ker(N^3)$};
  \node[layer] (dots) at (7.13,0) {$\cdots$};
  \node[layer] (nn) at (9.42,0) {$\ker(N^n)=\mathcal{V}$};

  \draw[arrow] (n1) -- (n2);
  \draw[arrow] (n2) -- (n3);
  \draw[arrow] (n3) -- (dots);
  \draw[arrow] (dots) -- (nn);
\end{tikzpicture}

  \caption{Extending bases through the kernel filtration makes $N$ move vectors to earlier layers.}
  \label{fig:nilpotent-filtration-basis}
\end{figure}

\begin{perspective} \label{per:generalized-eigenvectors-as-nilpotent-data} \hlt{Eigenvalue Plus Nilpotent Part}\\
On $\mathcal{G}_{\lambda}(T)$, the operator $T$ behaves as $\lambda I_{\mathcal{V}}$ plus the nilpotent operator $T-\lambda I_{\mathcal{V}}$. Generalized eigenvectors are therefore ordinary eigenvectors enriched by nilpotent chain data.
\end{perspective}

